% ------------------------------------------------------------------ Students' Submission
\chapter*{Students' Submission}
\addcontentsline{toc}{chapter}{Students' Submission}
\thispagestyle{plain}
We certify that this report is submitted as the FYP-1 mid-report for the project titled
\textit{\projecttitle}.

\vspace{1.2cm}
\begin{tabular}{p{7cm}l}
Student 1: Huzaifa Bin Khalid (22I-2332) & Signature: \rule{4.5cm}{0.4pt}\\[1.6cm]
Student 2: Taha Bashir Awan (23I-0824) & Signature: \rule{4.5cm}{0.4pt}\\[1.6cm]
Student 3: Saneed Ullah (23I-2568) & Signature: \rule{4.5cm}{0.4pt}\\
\end{tabular}

% ------------------------------------------------------------------ Anti-Plagiarism
\chapter*{Anti-Plagiarism Statement}
\addcontentsline{toc}{chapter}{Anti-Plagiarism Statement}
\thispagestyle{plain}
The students certify that this report is their own work and that all sources used in preparing
the report have been appropriately acknowledged and referenced.

% ------------------------------------------------------------------ Authors' Declaration
\chapter*{Authors' Declaration}
\addcontentsline{toc}{chapter}{Authors' Declaration}
\thispagestyle{plain}
We declare that the work documented in this report represents the work completed by the
project team up to the FYP-1 mid-report stage. Contributions, tools, external components,
and relevant sources are disclosed where applicable, including the use of an AI coding
assistant during platform implementation (Section~\ref{sec:tools-assisted}).

% ------------------------------------------------------------------ Abstract
\chapter*{Abstract}
\addcontentsline{toc}{chapter}{Abstract}
\thispagestyle{plain}
\textbf{Agricultural crop surveys still depend on field visits, manual crop identification and
hand-drawn GIS boundaries, which makes them slow, costly and hard to repeat over large regions.
GeoCrop is a satellite-to-GIS toolkit that learns crop classes from Sentinel-2 imagery and USDA
Cropland Data Layer (CDL) labels, maps cropland and crop types, derives field regions and area,
and will add yield estimation, all presented in a fast Web GIS portal. The research gap is that
existing tools and studies solve crop classification, field mapping or yield prediction in
isolation, and that research prototypes rarely remain usable once many images and model layers
must be displayed. The project follows an iterative plan with two parallel tracks. In Iteration 1
the architecture track replaced the proposal-stage prototype, which sent a 43.3~MB GeoJSON of
101,435 polygons to the browser, with a tile-based platform: Cloud Optimized GeoTIFF results served
by TiTiler, PMTiles vector layers, a React and MapLibre client, a FastAPI service, PostgreSQL with
PostGIS, Celery workers and S3-compatible storage, all run with Docker Compose. One screen of the
same area now costs 337~KB, cold map tiles return in 47~ms (median) and cached tiles in 1~ms, and an
uploaded Sentinel-2 image is processed by both crop models and published to the map in 57~s. A
manifest-based model registry lets new models plug in without changes to the API, database or
interface. The model track evaluated eleven crop-type models on two study areas with a spatial
block test; the proposed LightGBM pipeline with neighbourhood inputs and a 3$\times$3 neighbour vote
reaches 76.1\% accuracy (macro-F1 0.686) in eastern Washington and 92.2\% (macro-F1 0.803) in North
Dakota, while staying about 2~MB and CPU-only. Validation so far covers registry unit tests,
end-to-end API and browser tests, a rendering benchmark and the model evaluation protocol. FYP-2
will integrate the selected crop-type pipeline into the platform, add field-boundary and yield
models, run cross-area tests and prepare deployment.}

% ------------------------------------------------------------------ Acknowledgments
\chapter*{Acknowledgments}
\addcontentsline{toc}{chapter}{Acknowledgments}
\thispagestyle{plain}
The project team acknowledges the guidance and support received from the project supervisor,
Mr. Muhammad Aadil Ur Rehman, the external co-supervisor, Mr. Obaid Ur Rehman (Institute of Space
Technology), the evaluation panel whose proposal feedback shaped this iteration, and the faculty of
the FAST School of Computing. We also acknowledge the open data and open-source communities behind
Sentinel-2, the USDA Cropland Data Layer and the geospatial software used in this project.

% ------------------------------------------------------------------ Executive Summary
\chapter*{Executive Summary}
\addcontentsline{toc}{chapter}{Executive Summary}
\thispagestyle{plain}

\textbf{Project and motivation.} Agriculture departments need to know what is planted, where, over
how much area, and how much may be harvested. Today this comes from physical surveys followed by
manual GIS digitisation, which is time-consuming, resource-intensive and difficult to scale.
GeoCrop automates the chain from satellite imagery to crop maps, field regions, area statistics
and, in FYP-2, yield estimates, and presents them in a Web GIS portal (SDG~2, Zero Hunger).

\textbf{Scope.} The FYP builds a research-grade but usable prototype on open data: Sentinel-2
multispectral imagery as input, the USDA Cropland Data Layer (CDL) as training labels for eight
U.S. areas of interest (AOIs), and USDA NASS statistics for yield. Production-scale national
deployment, cadastral mapping and disease detection are out of scope.

\textbf{Research findings.} Two lines of investigation informed the design. On the platform side,
profiling the proposal prototype showed that its crop-type result for one AOI was a 43.3~MB GeoJSON
with 101,435 polygons (about 750,000 vertices), parsed three times and drawn as SVG, so loading and
zooming were slow and would not scale to more images or models. Keeping model output as rasters and
sending only the map tiles in view removes this cost. On the model side, eleven crop-type models
were compared on AOI1 (eastern Washington, 14 classes) and AOI2 (North Dakota, 9 classes) using
12 monthly Sentinel-2 composites of 2022 and a spatial block split. A time series raised AOI1
accuracy from 26.7\% (one date) to 75.2\%; the model family mattered little (eight models within
2 points); neighbourhood inputs and a 3$\times$3 vote added about 2 points; and 89\% of the remaining
errors sit on field-edge pixels. A small LightGBM model (about 2~MB, CPU-only) was proposed instead of
the much slower Presto foundation model. The proposal was approved with minor modifications; the
panel asked for clearer goals, evaluation criteria that measure performance, light weight, usability
and integration, and a feasibility analysis of the whole workflow. All three are addressed: goals
and criteria were restated with measurable targets, and every workflow stage was measured, showing
that the baseline crop-type model (3.1~GB RAM) is the bottleneck that LightGBM (549~MB) removes.

\textbf{Current implementation.} Iteration 1 delivered a working platform with seven containerised
services: a React and MapLibre Web GIS client, Nginx with a tile cache, a FastAPI service, a
TiTiler tile service, Celery workers, PostgreSQL with PostGIS and MinIO object storage. It supports
browsing and filtering layers, inspecting values at a point, field search, dashboard statistics,
uploading Sentinel-2 imagery, background model runs with live progress, a model catalogue and
system health. A model registry runs the current cropland and crop-type models and lets new models
plug in by folder and manifest.

\textbf{Validation status.} Registry unit tests pass (5/5), including a test that adds a model
without touching other code. End-to-end tests covered every API route, upload to map in 57~s,
crop filtering, field search and inspection in the browser. A benchmark measured 337~KB per screen
against 43.3~MB, cold tiles at 47~ms (median) and cached tiles at 1~ms. Crop-type accuracy is
reported against CDL labels on spatially separated test blocks, with confidence intervals.

\textbf{Risks.} CDL labels are themselves a model output; results cover one year and two areas; the
platform still runs the proposal-stage baseline model until the selected pipeline is integrated;
yield needs reliable historical ground truth; and boundary accuracy is limited at 30~m.

\textbf{Remaining work.} Iteration 2 (end of FYP-1) integrates the LightGBM crop-type pipeline,
adds the CDL confidence filter and a cross-area test, and supports multi-image upload. FYP-2 adds
field-boundary and yield models, an end-to-end evaluation against the success criteria, hardening
and deployment.

% ------------------------------------------------------------------ Abbreviations
\chapter*{List of Abbreviations}
\addcontentsline{toc}{chapter}{List of Abbreviations}
\thispagestyle{plain}
\begin{longtable}{@{}ll@{}}
AOI & Area of Interest\\
API & Application Programming Interface\\
CDL & Cropland Data Layer (USDA)\\
COG & Cloud Optimized GeoTIFF\\
CPU & Central Processing Unit\\
CRS & Coordinate Reference System\\
DEI & Diversity, Equity and Inclusion\\
ERD & Entity Relationship Diagram\\
F1 & Harmonic mean of precision and recall\\
FR & Functional Requirement\\
FYP & Final Year Project\\
GDAL & Geospatial Data Abstraction Library\\
GEE & Google Earth Engine\\
GIS & Geographic Information System\\
JWT & JSON Web Token\\
LightGBM & Light Gradient Boosting Machine\\
MLP & Multi-Layer Perceptron\\
MVT & Mapbox Vector Tile\\
NASS & National Agricultural Statistics Service (USDA)\\
NDVI & Normalized Difference Vegetation Index\\
NFR & Non-Functional Requirement\\
NIR & Near-Infrared\\
OSM & OpenStreetMap\\
PostGIS & Spatial extension for PostgreSQL\\
RAM & Random Access Memory\\
RF & Random Forest\\
S3 & Simple Storage Service (object storage API)\\
SDG & Sustainable Development Goal\\
SUS & System Usability Scale\\
SWIR & Short-Wave Infrared\\
TC & Test Case\\
UC & Use Case\\
USDA & United States Department of Agriculture\\
XYZ & Web map tile addressing scheme (zoom, column, row)\\
\end{longtable}
